\subsection*{Part 2, Draw shapes the short way: the shapes table}
Pentagon: $360 \div 5 = 72$. \quad Hexagon: $360 \div 6 = 60$. \quad Octagon: $360 \div 8 = 45$. \quad Decagon: $360 \div 10 = 36$. \quad Almost a circle: $360 \div 36 = 10$.

The 36-sided shape looks most like a circle. With \py|pencolor("red")| inside the loop, the whole triangle comes out red: line 4 makes the pen blue, but the red line runs before every \py|forward|, so blue is never used. With \py|print("Ta-da!")| at the margin it prints once; with four spaces it is inside the loop and prints 3 times.

\subsection*{Part 2, Stars: the star table}
\begin{center}
\renewcommand{\arraystretch}{1.3}
\begin{tabular}{P{2.4cm} P{3.4cm} P{2.4cm}}
\toprule
\textbf{Points} & \textbf{$180 \div$ points} & \textbf{Turn} \\
\midrule
9 & 20 & 160 \\
15 & 12 & 168 \\
45 & 4 & 176 \\
\bottomrule
\end{tabular}
\end{center}

\subsection*{Part 2, Square spiral: the missing line}
The first time round, \texttt{i} is 0, so \py|forward(i * 5)| is \py|forward(0)|. The turtle doesn't move, so the first ``line'' has no length. With \py|range(60)| and \py|forward(i * 3)|, the longest line is $59 \times 3 = 177$ steps, close to the $39 \times 5 = 195$ steps of the first spiral, so the size stays about the same while the lines sit closer together.

\subsection*{Part 2, Flower of squares}
36 squares $\times$ 10 degrees $= 360$, and 12 squares $\times$ 30 degrees $= 360$: the turns between the squares must add up to one full circle. \py|hideturtle()| needs \textbf{no} spaces, so it runs once, after the whole flower.

\subsection*{Part 2, Translate it yourself}
Your code may look a little different and still be right. If the turtle draws the same picture, you did it!
\begin{codenonum}
from ColabTurtlePlus.Turtle import *
clearscreen()
initializeTurtle((480, 360))
pensize(4)
pencolor("purple")
for i in range(8):
    for j in range(6):
        forward(40)
        right(60)
    right(45)
penup()
print("Done!")
\end{codenonum}
\py|forward(40)| and \py|right(60)| have 8 spaces (inside both loops). \py|right(45)| has 4 spaces. \py|penup()| and \py|print("Done!")| have none, because they come after the loops. You get 8 hexagons in a ring.

\subsection*{Part 2, Two-line puzzle}
\begin{codenonum}
for i in range(5):
    print(i + 1)
\end{codenonum}
\texttt{i} goes 0, 1, 2, 3, 4, so \py|i + 1| goes 1, 2, 3, 4, 5.
